% !TeX root = ../../Book2.tex
% 纲要标题：### 14.2 非交换 Nakayama 引理与生成元提升
% 纲要位置：计划纲要.md，第 1216 行。

\section{非交换 Nakayama 引理与生成元提升}
\label{b2:sec:1402}

% TODO: 按以下纲要要点撰写本节。

% 纲要内容（仅作写作计划，尚非教材正文）：
%
% * 对有限生成左 \(R\)-模 \(M\) 和双边理想 \(I\subseteq J(R)\)，证明 \(IM=M\Rightarrow M=0\)；用极小生成集与 \(1-a\) 的可逆性论证，不在非交换环上套用交换行列式技巧
% * 证明 \(M=N+IM\Rightarrow M=N\) 的生成元提升形式，并说明商模的有限生成性
% * 局部环和有限维代数中的剩余模、极小生成集及其维数解释；一般半单商下按简单模分量记录生成信息
% * 第三卷第3.2节保留交换环行列式技巧及局部版本的完整处理；第二卷第15章只调用本节已证的非交换形式
%

\subsection{非交换 Nakayama 引理}
\label{b2:subsec:1402:01}

待撰写。

\subsection{生成元提升与有限生成商模}
\label{b2:subsec:1402:02}

待撰写。

\subsection{剩余模与极小生成集}
\label{b2:subsec:1402:03}

待撰写。

\subsection{交换环版本与行列式技巧}
\label{b2:subsec:1402:04}

待撰写。
